\documentclass[10pt,a4paper]{article}

\usepackage{iftex}
\ifPDFTeX
  \usepackage[T2A]{fontenc}
  \usepackage[utf8]{inputenc}
  \usepackage[russian,english]{babel}
\else
  \usepackage{fontspec}
  \usepackage[russian,english]{babel}
  \setmainfont{DejaVu Serif}
  \setsansfont{DejaVu Sans}
  \setmonofont{DejaVu Sans Mono}
\fi
\usepackage{amsmath,amssymb}
\usepackage{array,longtable}
\usepackage{enumitem}
\usepackage{geometry}
\usepackage{hyperref}
\usepackage{url}

\geometry{left=18mm,right=18mm,top=16mm,bottom=18mm}
\setlength{\parindent}{3.5mm}
\setlength{\parskip}{0.25em}
\setlist{nosep,leftmargin=5mm}
\renewcommand{\arraystretch}{1.18}
\sloppy

\newcommand{\code}[1]{\texttt{#1}}
\newcommand{\A}[1]{\mathcal{A}_{#1}}
\newcommand{\Sync}{\operatorname{sync}}
\newcommand{\prio}{\operatorname{prio}}
\pdfstringdefDisableCommands{\def\A#1{A_#1}}

\title{Формализация PHY-уровня в иерархической SDN-управляемой 6G ISAC-сети}
\author{Редакция после формальной рецензии}
\date{6 июня 2026 г.}

\begin{document}

\maketitle

\section*{Аннотация}

В статье формализуется физический уровень PHY (Physical Layer) для иерархической модели SDN-управляемой сети 6G ISAC, где SDN (Software-Defined Networking) используется для оркестрации ресурсов, а ISAC (Integrated Sensing and Communication) объединяет передачу данных и радиозондирование среды. Модель не является радиофизическим симулятором и не утверждает, что обычные временные автоматы вычисляют SINR, CRB, вероятность обнаружения или распределения задержек. Измеренные PHY-величины и конфигурационные параметры сначала оцениваются внешним измерительным или аналитическим слоем, затем консервативно отображаются в конечные классы функцией абстракции \(\alpha_{PHY}:X_{meas}\times X_{cfg}\to X_{disc}\). Сеть PHY далее задается как композиция контрактных временных автоматов с явными часами, guards, resets, инвариантами, broadcast-уведомлениями отчетов, handshake-командами, приоритетами классификации, environment-stub автоматами и observer-автоматами для bounded deadlines. Такая постановка позволяет проверять не радиофизику как таковую, а своевременную реакцию MAC/SDN на классифицированные деградации communication и sensing KPI.

\section{Назначение и границы модели}

Физический уровень в 6G ISAC-сети выполняет две связанные функции: обеспечивает передачу данных и формирует информацию о физической среде. В отличие от детального моделирования формы сигнала на уровне I/Q-отсчетов, в работе используется абстрактная автоматная модель. Ее задача состоит в том, чтобы передавать вышестоящим уровням конечный набор KPI (Key Performance Indicators) и событий деградации, пригодных для model checking.

PHY-уровень не выбирает глобальную политику управления ресурсами. Выбор режима выполняется на уровнях MAC (Medium Access Control) и SDN. PHY применяет полученные конфигурации, классифицирует их последствия и формирует отчеты для контроллера. Поэтому статус модели следует понимать как \emph{абстрактную контрактную модель PHY для верификации реакций MAC/SDN на деградации ISAC-KPI}, а не как полную радиофизическую модель 6G-канала.

Обычные timed automata проверяют дискретную логику переходов во времени: достижимость, безопасность, отсутствие тупиков, соблюдение сроков и корректность синхронизации. Они не доказывают статистическое утверждение вида ``истинная вероятность обнаружения больше 0.95''. Вместо этого проверяется условие вида: если внешний PHY-estimator классифицировал \(Pd\) как \code{LOW}, то автомат \(\A{SQ}\) обязан за заданный срок перейти в состояние \code{ProbabilityLimited} и породить отчет \code{sensing\_report!}.

Литературная база в этой редакции оставлена точечно: signaling design для ISAC \cite{li2026rethinking}, SensCAP и cooperative sensing \cite{liu2024senscap,liu2026cooperative}, ISAC-enabled beam management \cite{huang2026beam}, теория timed automata \cite{alur1994timed}, семантика запросов UPPAAL \cite{uppaal2026queries} и пример применения формальной проверки к сетевой логике \cite{anand2007afdx}. Такой набор привязан к используемым утверждениям и не подменяет формальную часть библиографическим перечислением.

\section{Архитектурное положение PHY}

Рассматриваемая сеть задается иерархией:

\begin{verbatim}
Service/Application
       |
       v
SDN control plane
       |
       v
MAC/resource scheduling
       |
       v
PHY abstraction layer
       |
       v
Physical environment and external PHY estimators
\end{verbatim}

Сервисный уровень задает SLA: задержку, пропускную способность, надежность доставки, минимальную вероятность обнаружения, допустимую вероятность ложной тревоги, точность зондирования и допустимую давность информации. SDN-контроллер выбирает глобальную политику: разбиение ресурсов между связью и зондированием, маршрутизацию, восстановление и реконфигурацию. MAC-уровень реализует расписание радиоресурсов: плотность пилотов, распределение символов полезной нагрузки, мощность, MCS (Modulation and Coding Scheme) и команды формирования луча.

PHY является нижним измерительно-оценочным слоем. Он наблюдает радиосреду через внешние оценки, классифицирует состояние сигнала, канала, луча и sensing quality, после чего формирует события для MAC/SDN. Агрегирующий автомат \(\A{PH}\) получает не только отчет \(\A{SQ}\), но и отдельные отчеты \(\A{CH}\), \(\A{SIG}\) и \(\A{BM}\). Это устраняет архитектурную неоднозначность, при которой \(\A{PH}\) использовал communication degradation, но формально не имел входов от communication-компонентов.

\begin{verbatim}
Physical environment / estimators
        |             |             |
        v             v             v
      A_CH ------>  A_SIG ------>  A_SQ
        |             |             |
        v             v             v
      channel_     signal_      sensing_
      report!      report!      report!
        \             |             /
         \            |            /
          v           v           v
                    A_PH
                     |
                     v
              MAC and SDN control

A_BM receives beam commands from MAC/SDN
and publishes beam_report! to A_PH and A_SQ.
\end{verbatim}

Сигнальные технологии OFDM, OTFS и AFDM не задаются отдельными состояниями PHY. Они являются значениями дискретной переменной \(W\) в автомате \(\A{SIG}\). Состояния \(\A{SIG}\) описывают режим использования сигнала: pilot-based sensing, payload-assisted sensing, reconfiguration и degraded signal mode. Это предотвращает искусственное раздувание пространства состояний.

\section{Базовая модель контрактного timed automaton}

Каждый компонент PHY задается как контрактный временной автомат
\begin{equation}
\A{i}=(L_i,l_i^0,C_i,V_i,E_i,Inv_i,In_i,Out_i,Ass_i,Gar_i),
\end{equation}
где \(L_i\) --- конечное множество локаций, \(l_i^0\) --- начальная локация, \(C_i\) --- часы, \(V_i\) --- конечные переменные, \(E_i\) --- переходы, \(Inv_i\) --- инварианты локаций, \(In_i\) и \(Out_i\) --- входные и выходные каналы, \(Ass_i\) --- assumptions контракта, \(Gar_i\) --- guarantees контракта.

Переход имеет вид
\begin{equation}
e=(l,g,a,r,u,l'),
\end{equation}
где \(l,l'\in L_i\), \(g\) --- guard, \(a\in In_i\cup Out_i\cup\{\tau\}\) --- действие синхронизации или внутреннее действие, \(r\subseteq C_i\) --- множество сбрасываемых часов, \(u\) --- обновление конечных переменных. Переход разрешен, если истинны guard \(g\) и инвариант целевой локации после обновления \(u\) и сбросов \(r\).

Композиция самого PHY-уровня содержит ровно пять автоматов:
\begin{equation}
\A{PHY}=\A{CH}\parallel\A{SIG}\parallel\A{BM}\parallel\A{SQ}\parallel\A{PH}.
\end{equation}
Здесь нет шестого PHY-компонента: \(\A{ENV}\) не относится к PHY, не моделирует дополнительный сетевой слой и не принимает решений за MAC/SDN. Это служебная абстрактная среда, используемая только для закрытия модели при model checking.

Закрытая проверяемая система задается как композиция PHY-модели и environment-stub:
\begin{equation}
\A{SYS}=\A{PHY}\parallel\A{ENV},
\end{equation}
где
\begin{equation}
\A{ENV}=
\mathcal{A}_{\text{ENV\_CH}}\parallel
\mathcal{A}_{\text{ENV\_TARGET}}\parallel
\mathcal{A}_{\text{ENV\_MAC}}\parallel
\mathcal{A}_{\text{ENV\_NET}}.
\end{equation}
\(\A{ENV}\) кодирует assumptions о внешних входах: периодичность измерений, появление целей, допустимые команды MAC/SDN и доставку отчетов. Handshake-команды \code{x!}/\code{x?} синхронизируются бинарно по UPPAAL-подобной семантике, broadcast-отчеты и broadcast-события рассылаются всем готовым слушателям без перехвата, а внутренние \(\tau\)-переходы выполняются локально. Общие переменные записи отчетов обновляются только отправителем соответствующего отчета; остальные автоматы читают их после broadcast-уведомления.

\section{Абстракция непрерывной PHY-физики}

\subsection{Измерительный слой, конфигурация и конечные классы}

Непрерывные величины не вычисляются внутри timed automata. Они принадлежат внешнему слою измерений, аналитических оценок или симуляции. Для связи используются
\begin{equation}
SINR_c=\frac{S_c}{I_c+N_0B},
\end{equation}
где \(S_c\) --- полезная мощность communication-сигнала, \(I_c\) --- communication-интерференция, \(N_0B\) --- шумовая мощность в полосе \(B\).

Для зондирования используется отдельная величина
\begin{equation}
SINR_s=\frac{S_s}{I_{c\to s}+I_{self}+I_{mutual}+N_s},
\end{equation}
где \(S_s\) --- полезная sensing-компонента, \(I_{c\to s}\) --- интерференция communication-сигнала в sensing-тракте, \(I_{self}\) --- собственная интерференция, \(I_{mutual}\) --- взаимная интерференция между sensing-узлами, \(N_s\) --- шум sensing-приемника. Чтобы не смешивать физические измерения и дискретные настройки MAC/SDN, вход estimator-слоя разделяется на измерительный вектор и конфигурационный вектор:
\begin{equation}
\begin{aligned}
x_{PHY}^{meas}=(&P_t,P_r,S_c,S_s,N_0,N_s,I_c,I_{c\to s},I_{self},I_{mutual},L_p,\sigma_\tau,f_D,\\
&B,f_c,\Delta f,\epsilon_b,p_{mis},\Omega_{BM},Pd,Rfa,Acc_r,Acc_v,\\
&CRB_R,CRB_v,CRB_\theta,C_s,AoS_{BS},AoS_{CTRL},A_{cov}),\\
x_{PHY}^{cfg}=(&MCS,W,\rho_p,\kappa_{DRT},\psi_{cs},\psi_{ps},\theta_b,G_b,n_B,\\
&\tau_{SSB},T_{SSB},T_{PRS},N_{RE}^{PRS},\delta_{PRS},S_t,S_f,T_s,\gamma_{det}).
\end{aligned}
\end{equation}

Если часть величин не измеряется напрямую, она считается производной от измеренных параметров внешнего estimator-слоя. В любом случае конечная автоматная модель получает не вещественные значения из \(x_{PHY}^{meas}\) и не сырые настройки из \(x_{PHY}^{cfg}\), а только конечные классы из \(X_{disc}\). Поэтому guards автоматов далее задаются по class-переменным, а формулы с \(SINR\), \(Pd\), \(Rfa\), \(CRB\), \(\epsilon_b\) и \(\Omega_{BM}\) остаются на стороне estimator/score-слоя.

\subsection{Функция дискретизации}

Формальный мост между PHY-оценками, конфигурацией и timed automata задается функцией
\begin{equation}
\alpha_{PHY}:X_{meas}\times X_{cfg}\to X_{disc}.
\end{equation}
Ее результатом является набор bounded int или enum-переменных:
\begin{equation}
\begin{aligned}
X_{disc}=\{&SINRClass,BLERClass,CQIClass,IClass,DopplerClass,DelaySpreadClass,PowerClass,\\
&DRTClass,PilotDensityClass,PayloadSenseClass,PRSClass,BeamErrorClass,BlockageClass,\\
&PdClass,RfaClass,AccClass,CRBClass,AoSClass,CapClass,CoverageClass,\\
&ResourceShareClass,MisClass,BMOverheadClass,ChannelClass,SignalClass,BeamClass,\\
&SensingState,PHYState\}.
\end{aligned}
\end{equation}

Для communication SINR используется, например, следующая классификация:
\begin{equation}
SINRClass=\begin{cases}
\code{OUTAGE},& SINR_c<SINR_{out},\\
\code{LOW},& SINR_{out}\le SINR_c<SINR_{min},\\
\code{OK},& SINR_{min}\le SINR_c<SINR_{high},\\
\code{HIGH},& SINR_c\ge SINR_{high}.
\end{cases}
\end{equation}
Аналогично задаются классы:
\begin{align}
BLERClass&\in\{\code{OK},\code{HIGH},\code{CRITICAL}\},\\
CQIClass&\in\{\code{GOOD},\code{LIMITED},\code{BAD}\},\\
IClass&\in\{\code{OK},\code{HIGH},\code{CRITICAL}\},\\
DopplerClass&\in\{\code{OK},\code{HIGH}\},\\
DelaySpreadClass&\in\{\code{OK},\code{HIGH}\},\\
PowerClass&\in\{\code{OK},\code{LOW},\code{ABSENT}\},\\
DRTClass&\in\{\code{OK},\code{BAD}\},\\
PilotDensityClass&\in\{\code{OK},\code{LOW},\code{FAILED}\},\\
PayloadSenseClass&\in\{\code{OK},\code{LIMITED},\code{FAILED}\},\\
PRSClass&\in\{\code{OK},\code{LIMITED},\code{FAILED}\},\\
BeamErrorClass&\in\{\code{LOCKABLE},\code{UNSTABLE},\code{CRITICAL}\},\\
BlockageClass&\in\{\code{NONE},\code{SUSPECTED},\code{CONFIRMED}\},\\
PdClass&\in\{\code{OK},\code{LOW},\code{FAILED}\},\\
RfaClass&\in\{\code{OK},\code{HIGH},\code{CRITICAL}\},\\
AccClass&\in\{\code{OK},\code{LIMITED},\code{FAILED}\},\\
CRBClass&\in\{\code{OK},\code{LOOSE},\code{UNUSABLE}\},\\
AoSClass&\in\{\code{FRESH},\code{STALE},\code{EXPIRED}\},\\
CapClass&\in\{\code{OK},\code{LIMITED},\code{FAILED}\},\\
CoverageClass&\in\{\code{OK},\code{LIMITED},\code{FAILED}\},\\
ResourceShareClass&\in\{\code{OK},\code{LIMITED},\code{STARVED}\},\\
MisClass&\in\{\code{OK},\code{WARN},\code{CRITICAL}\},\\
BMOverheadClass&\in\{\code{OK},\code{HIGH},\code{EXCESSIVE}\},\\
ChannelClass&\in\{\code{NOMINAL},\code{INTERFERENCE\_LIMITED},\code{MOBILITY\_LIMITED},\\
&\quad\code{MULTIPATH\_LIMITED},\code{BLOCKAGE},\code{OUTAGE}\},\\
SignalClass&\in\{\code{NOMINAL},\code{PILOT\_BASED},\code{PAYLOAD\_ASSISTED},\\
&\quad\code{RECONFIGURING},\code{LIMITED}\},\\
BeamClass&\in\{\code{SEARCH},\code{TRACK},\code{LOCKED},\code{PREDICT},\\
&\quad\code{MISALIGNED},\code{RECOVERING},\code{HO\_ASSIST},\code{FAILED}\},\\
SensingState&\in\{\code{SensingQoSOk},\code{ProbabilityLimited},\code{FalseAlarmLimited},\\
&\quad\code{AccuracyLimited},\code{FreshnessLimited},\code{CoverageLimited},\\
&\quad\code{CapacityLimited},\code{SensingFailure}\},\\
PHYState&\in\{\code{PHYNormal},\code{PHYCommunicationDegraded},\\
&\quad\code{PHYSensingDegraded},\code{PHYJointDegraded},\\
&\quad\code{PHYRecovery},\code{PHYFailure}\}.
\end{align}

Для вероятностных и статистических KPI семантика следующая: \(Pd\), \(Rfa\), \(p_{mis}\) и CRB являются оценками estimator-слоя за окно наблюдения, а не вероятностями, доказанными внутри timed automata. Automata проверяют реакцию на классы этих оценок.

\subsection{Soundness и консервативность абстракции}

Абстракция выбирается как консервативная over-approximation относительно safety-свойств. Пограничные и неопределенные случаи относятся к худшему из соседних классов:
\begin{equation}
\alpha_{PHY}(x\in boundary(C_{good},C_{bad}))=C_{bad}.
\end{equation}
Например, при \(SINR_c=SINR_{min}\) и наличии неопределенности измерения \(\Delta SINR\) значение классифицируется как \code{LOW}, если нижняя граница доверительного интервала попадает ниже порога. Более строго это задается отношением конкретизации
\begin{equation}
R\subseteq (X_{meas}\times X_{cfg})\times X_{disc},
\end{equation}
где \((x,d)\in R\) означает, что дискретный класс \(d\) покрывает все допустимые estimator-значения и неопределенности для \(x\). Утверждение о сохранении safety-свойств далее не используется как доказанная теорема автоматически: оно предполагает корректную калибровку \(\alpha_{PHY}\), включение пограничных случаев в худшие классы и покрытие всех concrete transitions соответствующими abstract transitions относительно проверяемых atomic propositions. Если model checker находит counterexample, он может быть либо физически реализуемым сценарием, либо артефактом грубой over-approximation; поэтому counterexample должен проверяться обратным воспроизведением на estimator-слое или симуляторе.

\section{Семантика assume--guarantee контрактов}

Контракт автомата \(\A{i}\) трактуется не как описательная таблица, а как условное обязательство:
\begin{equation}
\A{i}\models Ass_i\Rightarrow Gar_i.
\end{equation}
Если assumptions выполняются, автомат обязан выполнить guarantees в указанных временных границах. Если assumption нарушается, автомат публикует соответствующее диагностическое событие из семейства \code{contract\_violation\_ch!}, \code{contract\_violation\_sig!}, \code{contract\_violation\_bm!}, \code{contract\_violation\_sq!}, \code{contract\_violation\_ph!} и переходит в диагностическую локацию \code{ContractViolation\_i}; guarantees после этого не используются как основание для доказательства свойств композиции.

Для исполнимой UPPAAL-подобной модели обозначения \(Ass_i\) и \(Gar_i\) не являются мета-переменными. Они реализуются как конкретные boolean-предикаты над clocks, class-переменными и состояниями environment-stub автоматов. Например, \code{Ass\_CH} означает, что \code{ENV\_CH} публикует \code{measure\_tick!} не реже \(T_{meas}+J_{meas}\), а все estimator-классы принадлежат объявленным enum-диапазонам; \code{Gar\_CH} означает, что после каждого принятого \code{measure\_tick?} до \(D_{meas}\) публикуется \code{channel\_report!} и обновляется \code{ChannelClass}. Аналогично \code{Ass\_SQ}, \code{Gar\_SQ}, \code{Ass\_PH} и \code{Gar\_PH} задаются как проверяемые условия, а не как свободные математические ярлыки.

Композиция корректна, если гарантии поставщика закрывают assumptions потребителя:
\begin{align}
Gar_{CH}&\Rightarrow Ass_{SQ},& Gar_{SIG}&\Rightarrow Ass_{SQ},& Gar_{BM}&\Rightarrow Ass_{SQ},\\
Gar_{CH}\wedge Gar_{SIG}\wedge Gar_{BM}\wedge Gar_{SQ}&\Rightarrow Ass_{PH},& Gar_{PH}&\Rightarrow Ass_{MAC/SDN}.
\end{align}

\begin{longtable}{|p{0.12\linewidth}|p{0.39\linewidth}|p{0.39\linewidth}|}
\hline
Автомат & Assumptions & Guarantees \\
\hline
\(\A{CH}\) & Measurements \(P_r,I_c,\sigma_\tau,f_D,SINR_c\) обновляются не реже \(T_{meas}\); estimator возвращает конечные классы. & За \(D_{meas}\) после \code{measure\_tick?} публикуется \code{channel\_report!}; \(ChannelClass\) выбирается по приоритету \(\prio_{CH}\). \\
\hline
\(\A{SIG}\) & MAC/SDN задают допустимые \(W,\rho_p,MCS,\kappa_{DRT},\psi_{cs},\psi_{ps}\), а estimator возвращает \code{PilotDensityClass}, \code{DRTClass} и \code{PayloadSenseClass}. & За \(D_{sig}\) после конфигурационной команды публикуется \code{signal\_report!}; нарушение DRT или BLER переводит автомат в \code{SignalLimited}. \\
\hline
\(\A{BM}\) & Команды луча и SSB/PRS-конфигурации конечны; \code{beam\_cmd?} доставляется за \(D_{cmd}\). & \code{beam\_report!} публикуется в каждом окне; misalignment порождает \code{beam\_misaligned!} за \(D_{BM}\); неуспешное recovery при \(c_{rec}=D_{BM}\) порождает \code{beam\_failure!}. \\
\hline
\(\A{SQ}\) & На вход поступают свежие \code{channel\_report?}, \code{signal\_report?} и \code{beam\_report?}. & За \(D_{sense}\) публикуется \code{sensing\_report!}; \(SensingState\) выбирается по приоритету \(\prio_{SQ}\). \\
\hline
\(\A{PH}\) & Все дочерние отчеты доставляются за \(D_{child}\); clocks \(aos_{bs}\), \(aos_{ctrl}\) корректно сбрасываются. & За \(D_{report}\) после деградации публикуется \code{phy\_kpi\_report!}; агрегированное состояние отражает communication и sensing classes. \\
\hline
\end{longtable}

\section{Operational semantics}

\subsection{Часы и resets}

Набор clocks делится на локальные clocks дочерних автоматов и observer clocks. Минимальный набор для проверяемой модели:
\begin{longtable}{|p{0.18\linewidth}|p{0.40\linewidth}|p{0.32\linewidth}|}
\hline
Clock & Смысл & Reset \\
\hline
\(c_{meas}\) & возраст последнего channel measurement & \code{measure\_tick?}, \code{channel\_report!} \\
\hline
\(c_{sig}\) & длительность применения signal configuration & \code{waveform\_config?}, \code{pilot\_config?}, \code{signal\_report!} \\
\hline
\(c_{sense}\) & время вычисления sensing state из входных классов & \code{channel\_report?}, \code{signal\_report?}, \code{beam\_report?}, \code{sensing\_report!} \\
\hline
\(c_{report}\) & возраст агрегированного PHY-отчета & \code{phy\_kpi\_report!} \\
\hline
\(c_{rec}\) & длительность beam recovery & вход в \code{BeamRecover} \\
\hline
\(c_{ssb}\) & длительность очередного SSB/beam-search окна & вход в \code{BeamSearch} или команда \code{ssb\_burst\_config?} \\
\hline
\(aos_{bs}\) & локальная давность sensing-информации на BS & \code{sensing\_success!} \\
\hline
\(aos_{ctrl}\) & давность sensing-информации у SDN-контроллера & доставка отчета контроллеру \\
\hline
\(c_{obs}\) & clock observer-автоматов для deadlines & вход observer в состояние ожидания \\
\hline
\end{longtable}

Основные инварианты:
\begin{align}
Inv(\code{MeasurePending})&: c_{meas}\le T_{meas}+J_{meas},\\
Inv(\code{SignalReconfiguring})&: c_{sig}\le D_{sig},\\
Inv(\code{SensingEvaluating})&: c_{sense}\le D_{sense},\\
Inv(\code{PHYKpiReporting})&: c_{report}\le D_{report},\\
Inv(\code{BeamRecover})&: c_{rec}\le D_{BM}.
\end{align}

\subsection{Синхронизационные каналы}

Каналы разделяются на бинарные handshake-команды среды и broadcast-уведомления отчетов/событий. Это важно для UPPAAL-семантики: один обычный handshake \code{x!} синхронизируется только с одним \code{x?}, тогда как отчеты \(\A{CH}\), \(\A{SIG}\), \(\A{BM}\), \(\A{SQ}\) и \(\A{PH}\) одновременно нужны нескольким потребителям и observer-автоматам. Поэтому report/event channels объявляются как \code{broadcast chan}; shared report variables обновляются отправителем перед публикацией, а получатели только читают их после broadcast-уведомления.
\begin{verbatim}
Handshake input from environment, MAC and SDN:
  chan measure_tick, pilot_config, prs_config, ssb_burst_config;
  chan waveform_config, payload_sensing_config, beam_cmd;
  chan extra_ssb_cmd, handover_assist_cmd, power_cmd;
  chan sensing_mode_cmd, recovery_cmd, new_beam_confirmed;
  chan mac_report_delivered, controller_report_delivered;

Broadcast report notifications:
  broadcast chan channel_report, signal_report, beam_report;
  broadcast chan sensing_report, phy_kpi_report;

Broadcast degradation, recovery and diagnostic events:
  broadcast chan channel_degraded, signal_degraded, beam_misaligned;
  broadcast chan beam_failure, sensing_degraded, phy_outage;
  broadcast chan degradation_event, recovery_event, mobility_alert;
  broadcast chan multipath_alert, blockage_detected, beam_locked;
  broadcast chan beam_restored, handover_hint, sensing_failure;
  broadcast chan phy_failure, sensing_success, aos_ctrl_expired;
  broadcast chan recovery_start, target_detected;
  broadcast chan contract_violation_ch, contract_violation_sig;
  broadcast chan contract_violation_bm, contract_violation_sq;
  broadcast chan contract_violation_ph;
\end{verbatim}

В таблицах ниже суффиксы \code{!}/\code{?} сохраняют направление использования канала. Для broadcast-канала запись \code{channel\_report!} означает публикацию уведомления, которое не блокируется отсутствием получателя; запись \code{channel\_report?} в \(\A{SQ}\), \(\A{PH}\) или observer означает пассивное получение того же уведомления без перехвата события у другого потребителя. Для handshake-команд, напротив, \code{recovery\_cmd!}/\code{recovery\_cmd?} выполняются парно между environment-stub и соответствующим автоматом. \(\A{SQ}\) использует broadcast-отчеты \(\A{CH}\), \(\A{SIG}\) и \(\A{BM}\), поскольку sensing quality зависит от канала, сигнальной конфигурации и качества сопровождения луча.

\subsection{Шаблон перехода}

Типовой переход report-автомата имеет одно действие синхронизации. Если нужны и событие деградации, и отчет, используются две последовательные дуги через committed/urgent report-pending локацию либо одно событие \code{report!}, в payload которого уже закодирован обновленный класс. В основной спецификации выбран второй вариант: degradation flags являются shared boolean/class fields внутри отчета, а отдельные события \code{channel\_degraded!}, \code{sensing\_degraded!} используются только там, где они показаны как отдельные переходы. Типовой переход имеет вид
\begin{equation}
(l,\;c\le D\wedge class=k,\;report!,\;\{c\},\;last\_class:=k,\;l').
\end{equation}
Если deadline истекает без публикации отчета, observer переходит в violation. Это важно: UPPAAL-запрос \(p\;--\!>\;q\) является unbounded leads-to и не задает дедлайн сам по себе. Bounded response задается отдельным clock или observer-автоматом.

\subsection{Абстрактная среда и закрытая модель}

Для проверки \code{A[] not deadlock}, reachability и deadline-свойств модель должна быть закрыта. Поэтому внешние входы задаются не неявно, а environment-stub автоматами с недетерминированным, но ограниченным поведением. Эти автоматы не входят в PHY-модель \(\A{PHY}\); они только формализуют окружение, с которым PHY взаимодействует:
\begin{verbatim}
ENV_CH:
  TickWait -- c_env <= T_meas + J_meas / measure_tick! --> TickWait
  update: estimator classes are refreshed before measure_tick!

ENV_TARGET:
  nondeterministically updates PRSClass, BeamErrorClass,
  BlockageClass and emits target_detected! when a target is visible.

ENV_MAC:
  emits pilot_config!, waveform_config!, recovery_cmd!,
  extra_ssb_cmd! and new_beam_confirmed! under admissible configs.

ENV_NET:
  delivers mac_report_delivered! and controller_report_delivered!
  within bounded network delay D_net.
\end{verbatim}
Если эти автоматы не включены в композицию, спецификация трактуется как open-system contract, а свойства формулируются условно: они должны выполняться только при \code{Ass\_ENV}, то есть при соблюдении ограничений на частоту измерений, доставку команд, обновление estimator-классов и сетевые задержки.

\section{Автоматы PHY}

\subsection{Автомат канала A\_CH}

\(\A{CH}\) классифицирует состояние канала по дискретным классам, а не по конкурирующим вещественным условиям. Чтобы избежать недетерминизма при одновременном выполнении нескольких условий, вводится приоритет:
\begin{equation}
\begin{aligned}
\prio_{CH}:\quad &\code{Outage}>\code{Blockage}>\code{InterferenceLimited}>\\
&\code{MobilityLimited}>\code{MultipathLimited}>\code{ChannelNominal}.
\end{aligned}
\end{equation}
Функция классификации задается явно:
\begin{equation}
ChannelClass=highest\_priority\_CH(SINRClass,PowerClass,IClass,DopplerClass,DelaySpreadClass),
\end{equation}
где \code{OUTAGE} выбирается при \(SINRClass=\code{OUTAGE}\), \code{BLOCKAGE} --- при \(PowerClass\in\{\code{LOW},\code{ABSENT}\}\) без более приоритетного outage, \code{INTERFERENCE\_LIMITED} --- при \(IClass\in\{\code{HIGH},\code{CRITICAL}\}\) без outage/blockage, далее по \(DopplerClass\) и \(DelaySpreadClass\). Если estimator одновременно возвращает \(IClass=\code{CRITICAL}\), \(PowerClass=\code{LOW}\) и \(SINRClass=\code{OUTAGE}\), выбирается \code{OUTAGE}; вещественные \(I_c\) и \(P_r\) в guards автомата не используются.

\begin{longtable}{|p{0.24\linewidth}|p{0.35\linewidth}|p{0.30\linewidth}|}
\hline
Локация & Guard с учетом \(\prio_{CH}\) & Output / update \\
\hline
\code{ChannelNominal} & \(SINRClass\in\{\code{OK},\code{HIGH}\}\) и нет более приоритетных нарушений & \code{channel\_report!}, \(ChannelClass:=\code{NOMINAL}\) \\
\hline
\code{InterferenceLimited} & \(IClass\in\{\code{HIGH},\code{CRITICAL}\}\) и нет \code{Outage}, \code{Blockage} & \code{channel\_report!}, \(ChannelClass:=\code{INTERFERENCE\_LIMITED}\) \\
\hline
\code{MobilityLimited} & \(DopplerClass=\code{HIGH}\) и нет более приоритетных нарушений & \code{channel\_report!}, \(ChannelClass:=\code{MOBILITY\_LIMITED}\) \\
\hline
\code{MultipathLimited} & \(DelaySpreadClass=\code{HIGH}\) и нет более приоритетных нарушений & \code{channel\_report!}, \(ChannelClass:=\code{MULTIPATH\_LIMITED}\) \\
\hline
\code{Blockage} & \(PowerClass\in\{\code{LOW},\code{ABSENT}\}\) и \(SINRClass\ne\code{OUTAGE}\) & \code{channel\_report!}, \(ChannelClass:=\code{BLOCKAGE}\) \\
\hline
\code{Outage} & \(SINRClass=\code{OUTAGE}\) & \code{channel\_report!}, \(ChannelClass:=\code{OUTAGE}\) \\
\hline
\end{longtable}

Явная схема переходов \(E_{CH}\) отделяет прием измерения от публикации отчета:
\begin{verbatim}
Any -- measure_tick? / c_meas := 0 --> MeasurePending
MeasurePending -- guard: ChannelClass' == NOMINAL
  sync: channel_report!
  update: ChannelClass := NOMINAL, channel_degraded_flag := false,
          c_meas := 0 --> ChannelNominal
MeasurePending -- guard: ChannelClass' == INTERFERENCE_LIMITED
  sync: channel_report!
  update: ChannelClass := INTERFERENCE_LIMITED,
          channel_degraded_flag := true, c_meas := 0 --> InterferenceLimited
MeasurePending -- guard: ChannelClass' == MOBILITY_LIMITED
  sync: channel_report!
  update: ChannelClass := MOBILITY_LIMITED,
          channel_degraded_flag := true, c_meas := 0 --> MobilityLimited
MeasurePending -- guard: ChannelClass' == MULTIPATH_LIMITED
  sync: channel_report!
  update: ChannelClass := MULTIPATH_LIMITED,
          channel_degraded_flag := true, c_meas := 0 --> MultipathLimited
MeasurePending -- guard: ChannelClass' == BLOCKAGE
  sync: channel_report!
  update: ChannelClass := BLOCKAGE,
          channel_degraded_flag := true, c_meas := 0 --> Blockage
MeasurePending -- guard: ChannelClass' == OUTAGE
  sync: channel_report!
  update: ChannelClass := OUTAGE,
          channel_degraded_flag := true, c_meas := 0 --> Outage
\end{verbatim}
Здесь \code{ChannelClass'} обозначает результат функции \code{highest\_priority\_CH} на свежих class-входах. Если требуется отдельный observable alert, он отправляется отдельной последующей broadcast-дугой из соответствующей локации, например \code{Outage -- phy\_outage! --> Outage}; она не совмещается с \code{channel\_report!} на одном edge.

\subsection{Автомат сигнальной конфигурации A\_SIG}

Переменная waveform остается конечной переменной
\begin{equation}
W\in\{\code{OFDM},\code{OTFS},\code{AFDM},\code{SC},\code{OTHER}\}.
\end{equation}
Локации \(\A{SIG}\) описывают не waveform, а режим эксплуатации сигнала.

\begin{longtable}{|p{0.25\linewidth}|p{0.36\linewidth}|p{0.29\linewidth}|}
\hline
Локация & Guard / invariant & Output / transition \\
\hline
\code{SignalNominal} & \(BLERClass=\code{OK}\), \(DRTClass=\code{OK}\) & \code{signal\_report!}, \(SignalClass:=\code{NOMINAL}\); при \(PilotDensityClass=\code{OK}\) возможен переход в \code{PilotBasedSensing} \\
\hline
\code{PilotBasedSensing} & \(PilotDensityClass=\code{OK}\); \(c_{sig}\le D_{sig}\) & \code{signal\_report!}, \(SignalClass:=\code{PILOT\_BASED}\); при \(PilotDensityClass=\code{LOW}\) и \(PayloadSenseClass\ne\code{FAILED}\) переход в \code{PayloadAssistedSensing} \\
\hline
\code{Payload}\-\code{Assisted}\-\code{Sensing} & \(PayloadSenseClass\ne\code{FAILED}\) & \code{signal\_report!}, \(SignalClass:=\code{PAYLOAD\_ASSISTED}\); при \(DRTClass=\code{BAD}\) или \(BLERClass\ne\code{OK}\) переход в \code{SignalLimited} \\
\hline
\code{SignalReconfiguring} & получена \code{waveform\_config?} или \code{pilot\_config?}; \(c_{sig}\le D_{sig}\) & \code{signal\_report!}, \(SignalClass:=\code{RECONFIGURING}\), \(c_{sig}:=0\) \\
\hline
\code{SignalLimited} & \(BLERClass\in\{\code{HIGH},\code{CRITICAL}\}\) или \(DRTClass=\code{BAD}\) или \(PayloadSenseClass=\code{FAILED}\) & \code{signal\_report!}, \(SignalClass:=\code{LIMITED}\); ожидание реконфигурации MAC/SDN \\
\hline
\end{longtable}

Явная схема переходов \(E_{SIG}\):
\begin{verbatim}
Any -- waveform_config? / c_sig := 0 --> SignalReconfiguring
Any -- pilot_config? / c_sig := 0 --> SignalReconfiguring
Any -- payload_sensing_config? / c_sig := 0 --> SignalReconfiguring
SignalReconfiguring -- c_sig <= D_sig
  sync: signal_report!
  update: SignalClass := RECONFIGURING, c_sig := 0 --> SignalNominal
SignalNominal -- PilotDensityClass == OK && BLERClass == OK && DRTClass == OK
  sync: signal_report!
  update: SignalClass := PILOT_BASED, c_sig := 0 --> PilotBasedSensing
PilotBasedSensing -- PilotDensityClass == LOW && PayloadSenseClass != FAILED
  sync: signal_report!
  update: SignalClass := PAYLOAD_ASSISTED, c_sig := 0 --> PayloadAssistedSensing
Any -- BLERClass in {HIGH, CRITICAL} || DRTClass == BAD
       || PayloadSenseClass == FAILED
  sync: signal_report!
  update: SignalClass := LIMITED, signal_degraded_flag := true,
          c_sig := 0 --> SignalLimited
SignalLimited -- waveform_config? / c_sig := 0 --> SignalReconfiguring
\end{verbatim}
Отдельное событие \code{signal\_degraded!} при необходимости выводится отдельной broadcast-дугой из \code{SignalLimited}; оно не совмещается с \code{signal\_report!} на том же переходе.

\subsection{Автомат beam management A\_BM}

\(\A{BM}\) описывает начальный поиск направления, сопровождение цели, фиксацию узкого луча, прогноз выхода цели за границу луча, восстановление направления и помощь handover. Он не вычисляет оптимальный beamforming vector; соответствующая оптимизация остается вне timed automata. Автомат получает конечные классы \(MisClass\), \(BlockageClass\), \(PRSClass\) и \(BMOverheadClass\), публикуя события для MAC/SDN.

Классификация рассогласования задается как
\begin{equation}
MisClass=\begin{cases}
\code{OK},& p_{mis}\le p_{mis}^{warn},\\
\code{WARN},& p_{mis}^{warn}<p_{mis}\le p_{mis}^{max},\\
\code{CRITICAL},& p_{mis}>p_{mis}^{max}.
\end{cases}
\end{equation}
Класс накладных расходов получается из
\begin{equation}
\Omega_{BM}=\min\left(1,\frac{N_{SSB}T_{SSB}+N_{PRS}T_{PRS}+T_{ctrl}+T_{report}}{T_{frame}}\right).
\end{equation}

Ключевая правка касается recovery-deadline. В локации \code{BeamRecover} действует инвариант \(c_{rec}\le D_{BM}\). Поэтому переход в \code{BeamFailed} не может иметь guard \(c_{rec}>D_{BM}\), так как такое состояние времени недостижимо. Корректный guard задается только через clock и классы:
\begin{equation}
\code{BeamRecover}\xrightarrow{c_{rec}=D_{BM}\wedge BeamErrorClass\ne\code{LOCKABLE}\wedge BlockageClass\ne\code{CONFIRMED}\;/\;beam\_failure!}\code{BeamFailed}.
\end{equation}

\begin{longtable}{|p{0.18\linewidth}|p{0.42\linewidth}|p{0.30\linewidth}|}
\hline
Локация & Guard / invariant & Output / transition \\
\hline
\code{BeamSearch} & широкое сканирование SSB; \(c_{ssb}\le \tau_{SSB}+J_{SSB}\) & при \code{target\_detected?} и \(PRSClass\ne\code{FAILED}\): \(BeamClass:=\code{TRACK}\), переход в \code{BeamTrack} \\
\hline
\code{BeamTrack} & PRS-оценка положения и скорости & при \(BeamErrorClass=\code{LOCKABLE}\): \(BeamClass:=\code{LOCKED}\), переход в \code{BeamLock}; при \(AccClass=\code{FAILED}\) или \(BeamErrorClass=\code{CRITICAL}\): \(BeamClass:=\code{MISALIGNED}\) \\
\hline
\code{BeamLock} & \(MisClass=\code{OK}\) & \code{beam\_report!}, \(BeamClass:=\code{LOCKED}\); при \(MisClass=\code{WARN}\) или \(BlockageClass=\code{SUSPECTED}\) переход в \code{BeamPredict} \\
\hline
\code{BeamPredict} & прогноз выхода цели за границу луча или блокировки & при \code{extra\_ssb\_cmd?} переход в \code{BeamRecover}, \(c_{rec}:=0\) \\
\hline
\code{BeamMisalign} & \(MisClass=\code{CRITICAL}\) или \(BeamErrorClass=\code{CRITICAL}\) & \code{beam\_misaligned!}; затем переход в \code{BeamRecover}, \(c_{rec}:=0\), \(BeamClass:=\code{RECOVERING}\) \\
\hline
\code{BeamRecover} & \(c_{rec}\le D_{BM}\) & success: \(BeamErrorClass=\code{LOCKABLE}\), \code{beam\_restored!}, \code{BeamLock}; blockage: \(BlockageClass=\code{CONFIRMED}\), \code{handover\_hint!}, \code{BeamHOAssist}; timeout: \(c_{rec}=D_{BM}\), \code{beam\_failure!} \\
\hline
\code{BeamHOAssist} & требуется помощь handover & при \code{new\_beam\_confirmed?} переход в \code{BeamTrack} \\
\hline
\code{BeamFailed} & recovery не завершен в срок & ожидание \code{recovery\_cmd?} или новой MAC/SDN-конфигурации \\
\hline
\end{longtable}

Явная схема переходов \(E_{BM}\) использует один sync на каждой дуге:
\begin{verbatim}
Any -- beam_cmd? / BeamClass := SEARCH, c_ssb := 0 --> BeamSearch
BeamSearch -- target_detected? / target_seen := true --> BeamSearchSeen
BeamSearchSeen -- PRSClass != FAILED
  sync: beam_report!
  update: BeamClass := TRACK --> BeamTrack
BeamTrack -- BeamErrorClass == LOCKABLE
  sync: beam_report!
  update: BeamClass := LOCKED --> BeamLock
BeamTrack -- AccClass == FAILED || BeamErrorClass == CRITICAL
  sync: beam_report!
  update: BeamClass := MISALIGNED, beam_degraded_flag := true --> BeamMisalign
BeamLock -- MisClass == OK
  sync: beam_report!
  update: BeamClass := LOCKED --> BeamLock
BeamLock -- MisClass == WARN || BlockageClass == SUSPECTED
  sync: beam_report!
  update: BeamClass := PREDICT --> BeamPredict
BeamPredict -- extra_ssb_cmd? / c_rec := 0 --> BeamRecoveryStart
BeamRecoveryStart -- recovery_start!
  update: BeamClass := RECOVERING --> BeamRecover
BeamMisalign -- beam_misaligned! --> BeamRecoveryStart
BeamRecover -- BeamErrorClass == LOCKABLE && c_rec <= D_BM
  sync: beam_restored!
  update: BeamClass := LOCKED --> BeamLock
BeamRecover -- BlockageClass == CONFIRMED && c_rec <= D_BM
  sync: handover_hint!
  update: BeamClass := HO_ASSIST --> BeamHOAssist
BeamRecover -- c_rec == D_BM && BeamErrorClass != LOCKABLE
              && BlockageClass != CONFIRMED
  sync: beam_failure!
  update: BeamClass := FAILED --> BeamFailed
BeamHOAssist -- new_beam_confirmed? / BeamClass := TRACK --> BeamTrack
BeamFailed -- recovery_cmd? / BeamClass := RECOVERING,
              c_rec := 0 --> BeamRecover
\end{verbatim}
Периодический \code{beam\_report!} может публиковаться из \code{BeamRecover}, \code{BeamHOAssist} и \code{BeamFailed} отдельными self-loop переходами с единственным sync \code{beam\_report!}; observer recovery слушает outcome-события \code{beam\_restored!}, \code{handover\_hint!} и \code{beam\_failure!}, а не перехватывает \code{beam\_report!}.

Контракт \(\A{BM}\) в проверяемой форме:
\begin{verbatim}
Assume:
  n_B is finite;
  SSB/PRS configuration is admissible;
  beam_cmd? is delivered within D_cmd.

Guarantee:
  if MisClass = CRITICAL then beam_misaligned! is emitted within D_BM;
  if BeamRecover is entered then exactly one of
      beam_restored!, handover_hint!, beam_failure!
    is emitted not later than c_rec = D_BM;
  Omega_BM and BMOverheadClass are included in every beam_report!.
\end{verbatim}

\subsection{Автомат качества зондирования A\_SQ}

\(\A{SQ}\) получает \code{channel\_report?}, \code{signal\_report?} и \code{beam\_report?}, затем вычисляет конечное состояние sensing quality. Все guards задаются по классам, а не по вещественным \(Pd\), \(Rfa\), \(Acc_r\), \(Acc_v\), CRB или \(C_s\).

\begin{equation}
\begin{aligned}
\prio_{SQ}:\quad &\code{SensingFailure}>\code{FreshnessLimited}>\code{AccuracyLimited}>\\
&\code{ProbabilityLimited}>\code{FalseAlarmLimited}>\code{CapacityLimited}>\\
&\code{CoverageLimited}>\code{SensingQoSOk}.
\end{aligned}
\end{equation}

\begin{longtable}{|p{0.22\linewidth}|p{0.42\linewidth}|p{0.26\linewidth}|}
\hline
Локация & Guard по классам & Output / update \\
\hline
\code{SensingQoSOk} & все sensing-классы равны \code{OK}/\code{FRESH} и нет более приоритетных нарушений & \code{sensing\_report!}, \(SensingState:=\code{SensingQoSOk}\) \\
\hline
\code{ProbabilityLimited} & \(PdClass=\code{LOW}\) и нет более приоритетных нарушений & \code{sensing\_report!}, \(SensingState:=\code{ProbabilityLimited}\) \\
\hline
\code{FalseAlarmLimited} & \(RfaClass=\code{HIGH}\) и нет более приоритетных нарушений & \code{sensing\_report!}, \(SensingState:=\code{FalseAlarmLimited}\) \\
\hline
\code{AccuracyLimited} & \(AccClass=\code{LIMITED}\) или \(CRBClass=\code{LOOSE}\) & \code{sensing\_report!}, \(SensingState:=\code{AccuracyLimited}\) \\
\hline
\code{FreshnessLimited} & \(AoSClass=\code{STALE}\) и нет \code{SensingFailure} & \code{sensing\_report!}, \(SensingState:=\code{FreshnessLimited}\) \\
\hline
\code{CoverageLimited} & \(CoverageClass=\code{LIMITED}\) и нет более приоритетных нарушений & \code{sensing\_report!}, \(SensingState:=\code{CoverageLimited}\) \\
\hline
\code{CapacityLimited} & \(CapClass=\code{LIMITED}\) и нет более приоритетных нарушений & \code{sensing\_report!}, \(SensingState:=\code{CapacityLimited}\) \\
\hline
\code{SensingFailure} & \(PdClass=\code{FAILED}\) или \(RfaClass=\code{CRITICAL}\) или \(AccClass=\code{FAILED}\) или \(CRBClass=\code{UNUSABLE}\) или \(AoSClass=\code{EXPIRED}\) или \(CapClass=\code{FAILED}\) или \(CoverageClass=\code{FAILED}\) или \(BeamClass=\code{FAILED}\) & \code{sensing\_report!}, \(SensingState:=\code{SensingFailure}\) \\
\hline
\end{longtable}

Детерминированность классификации проверяется счетчиком enabled-guards или явной функцией:
\begin{equation}
\begin{aligned}
SensingState={}&highest\_priority\_SQ(PdClass,RfaClass,AccClass,CRBClass,\\
&AoSClass,CapClass,CoverageClass,BeamClass).
\end{aligned}
\end{equation}

Явная схема переходов \(E_{SQ}\):
\begin{verbatim}
Idle -- channel_report? / input_changed := true, c_sense := 0 --> SensingEvaluating
Idle -- signal_report?  / input_changed := true, c_sense := 0 --> SensingEvaluating
Idle -- beam_report?    / input_changed := true, c_sense := 0 --> SensingEvaluating
SensingEvaluating -- guard: SensingState' == SensingQoSOk
  sync: sensing_report!
  update: SensingState := SensingQoSOk,
          sensing_degraded_flag := false, c_sense := 0 --> SensingQoSOk
SensingEvaluating -- guard: SensingState' == ProbabilityLimited
  sync: sensing_report!
  update: SensingState := ProbabilityLimited,
          sensing_degraded_flag := true, c_sense := 0 --> ProbabilityLimited
SensingEvaluating -- guard: SensingState' == FalseAlarmLimited
  sync: sensing_report!
  update: SensingState := FalseAlarmLimited,
          sensing_degraded_flag := true, c_sense := 0 --> FalseAlarmLimited
SensingEvaluating -- guard: SensingState' == AccuracyLimited
  sync: sensing_report!
  update: SensingState := AccuracyLimited,
          sensing_degraded_flag := true, c_sense := 0 --> AccuracyLimited
SensingEvaluating -- guard: SensingState' == FreshnessLimited
  sync: sensing_report!
  update: SensingState := FreshnessLimited,
          sensing_degraded_flag := true, c_sense := 0 --> FreshnessLimited
SensingEvaluating -- guard: SensingState' == CoverageLimited
  sync: sensing_report!
  update: SensingState := CoverageLimited,
          sensing_degraded_flag := true, c_sense := 0 --> CoverageLimited
SensingEvaluating -- guard: SensingState' == CapacityLimited
  sync: sensing_report!
  update: SensingState := CapacityLimited,
          sensing_degraded_flag := true, c_sense := 0 --> CapacityLimited
SensingEvaluating -- guard: SensingState' == SensingFailure
  sync: sensing_report!
  update: SensingState := SensingFailure,
          sensing_failure_flag := true, c_sense := 0 --> SensingFailure
\end{verbatim}
Здесь \code{SensingState'} --- результат функции \code{highest\_priority\_SQ}. Broadcast-события \code{sensing\_degraded!} и \code{sensing\_failure!} могут публиковаться отдельными self-loop переходами после обновления флагов; основной отчет остается единственным sync на report-edge.

\subsection{Агрегирующий автомат A\_PH}

\(\A{PH}\) поддерживает две независимые булевы переменные:
\begin{align}
comm\_ok &\equiv (ChannelClass\in\{\code{NOMINAL},\code{INTERFERENCE\_LIMITED}\}\vee{}\\
&\quad ChannelClass\in\{\code{MOBILITY\_LIMITED},\code{MULTIPATH\_LIMITED}\})\wedge{}\\
&\quad SignalClass\ne\code{LIMITED}\wedge BeamClass\notin\{\code{FAILED},\code{MISALIGNED}\},\\
sensing\_qos\_ok &\equiv SensingState=\code{SensingQoSOk}.
\end{align}
Эти переменные обновляются только после broadcast-уведомлений \code{channel\_report?}, \code{signal\_report?}, \code{beam\_report?} и \code{sensing\_report?}. Агрегированная переменная \code{PHYState} обновляется только \(\A{PH}\), чтобы остальные автоматы читали единый источник истины.

\begin{longtable}{|p{0.31\linewidth}|p{0.34\linewidth}|p{0.25\linewidth}|}
\hline
Локация & Guard & Output / transition \\
\hline
\code{PHYNormal} & \(comm\_ok\wedge sensing\_qos\_ok\) & \code{phy\_kpi\_report!}, \(PHYState:=\code{PHYNormal}\) \\
\hline
\code{PHYCommunicationDegraded} & \(\neg comm\_ok\wedge sensing\_qos\_ok\) & \code{phy\_kpi\_report!}, \(PHYState:=\code{PHYCommunicationDegraded}\) \\
\hline
\code{PHYSensingDegraded} & \(comm\_ok\wedge\neg sensing\_qos\_ok\) & \code{phy\_kpi\_report!}, \(PHYState:=\code{PHYSensingDegraded}\) \\
\hline
\code{PHYJointDegraded} & \(\neg comm\_ok\wedge\neg sensing\_qos\_ok\) & \code{phy\_kpi\_report!}, \(PHYState:=\code{PHYJointDegraded}\) \\
\hline
\code{PHYKpiReporting} & истек \(T_{report}\) или получено событие деградации; \(c_{report}\le D_{report}\) & публикация \code{phy\_kpi\_report!}, \(c_{report}:=0\) \\
\hline
\code{PHYRecovery} & получена \code{recovery\_cmd?} & \(PHYState:=\code{PHYRecovery}\); при восстановлении классов переход в \code{PHYNormal}; иначе deadline observer фиксирует нарушение \\
\hline
\code{PHYFailure} & критический отказ канала, луча или sensing-функции & \code{phy\_failure!}, \(PHYState:=\code{PHYFailure}\); ожидание внешней реконфигурации \\
\hline
\end{longtable}

Явная схема переходов \(E_{PH}\):
\begin{verbatim}
Any -- channel_report? / child_update := true, c_report := 0 --> PHYKpiReporting
Any -- signal_report?  / child_update := true, c_report := 0 --> PHYKpiReporting
Any -- beam_report?    / child_update := true, c_report := 0 --> PHYKpiReporting
Any -- sensing_report? / child_update := true, c_report := 0 --> PHYKpiReporting
PHYKpiReporting -- comm_ok && sensing_qos_ok && c_report <= D_report
  sync: phy_kpi_report!
  update: PHYState := PHYNormal, c_report := 0 --> PHYNormal
PHYKpiReporting -- !comm_ok && sensing_qos_ok && c_report <= D_report
  sync: phy_kpi_report!
  update: PHYState := PHYCommunicationDegraded,
          degradation_flag := true, c_report := 0 --> PHYCommunicationDegraded
PHYKpiReporting -- comm_ok && !sensing_qos_ok && c_report <= D_report
  sync: phy_kpi_report!
  update: PHYState := PHYSensingDegraded,
          degradation_flag := true, c_report := 0 --> PHYSensingDegraded
PHYKpiReporting -- !comm_ok && !sensing_qos_ok && c_report <= D_report
  sync: phy_kpi_report!
  update: PHYState := PHYJointDegraded,
          degradation_flag := true, c_report := 0 --> PHYJointDegraded
Any -- recovery_cmd? / PHYState := PHYRecovery --> PHYRecovery
Any -- ChannelClass == OUTAGE || BeamClass == FAILED
       || SensingState == SensingFailure
  sync: phy_failure!
  update: PHYState := PHYFailure --> PHYFailure
\end{verbatim}
Если нужен отдельный \code{degradation\_event!}, он публикуется дополнительным broadcast self-loop после установки \code{degradation\_flag}; основной KPI-отчет остается единственным sync на переходе \(\A{PH}\).

\section{SensCAP, capacity и freshness}

SensCAP трактуется как системная метрика, включающая sensing probability, sensing accuracy и sensing capacity \cite{liu2024senscap,liu2026cooperative}. В автоматной модели эти величины не используются как непрерывные guards напрямую; они предварительно переводятся в классы.

Sensing probability задается estimator-слоем как
\begin{equation}
Pd(R_{fa})=\frac{N_{det}}{N},\qquad R_{fa}\le R_{fa}^{max}.
\end{equation}
Для CFAR-приближения порог может быть записан как
\begin{equation}
S_{thres}=N_{ru}\left((R_{fa})^{-1/N_{ru}}-1\right)p_n.
\end{equation}
Точность зондирования задается через условия
\begin{equation}
\Pr(\Delta r\le l_a)\ge q,\qquad \Pr(\Delta v\le v_a)\ge q_v,
\end{equation}
а CRB сохраняется как теоретическая нижняя граница ошибки
\begin{equation}
\mathbf{CRB}=(CRB_R,CRB_v,CRB_\theta).
\end{equation}

Формула sensing capacity должна явно содержать временное окно и требования качества:
\begin{equation}
C_s(Q_s,T_0)=\frac{N^*(Q_s,T_0)}{A},
\end{equation}
где \(A\) --- площадь зоны зондирования, \(T_0\) --- окно оценки, \(Q_s=(l_a,v_a,Pd_{min},Rfa_{max},q,q_v,T_0)\), а \(N^*(Q_s,T_0)\) --- максимальное число целей, для которого выполняются требования
\begin{align}
Pd(R_{fa}^{max})&\ge Pd_{min},\\
\Pr(\Delta r\le l_a)&\ge q,\\
\Pr(\Delta v\le v_a)&\ge q_v.
\end{align}
Доля ресурсов зондирования задается как
\begin{equation}
u_s=\frac{T_s}{T_0},\qquad 0<u_s\le 1.
\end{equation}
В timed automata используются только \(CapClass\) и \(ResourceShareClass\), например \code{OK}, \code{LIMITED}, \code{FAILED}.

AoS (Age of Sensing) разделяется на локальную и контроллерную метрики:
\begin{align}
AoS_{BS}(t)&=t-t_{sense}^{last},\\
AoS_{CTRL}(t)&=t-t_{ctrl}^{last}.
\end{align}
\(AoS_{BS}\) сбрасывается после \code{sensing\_success!}, а \(AoS_{CTRL}\) --- после \code{controller\_report\_delivered?}. Поэтому \(AoS_{CTRL}\) может быть больше \(AoS_{BS}\) при перегрузке, изменении маршрута или задержке обработки на edge/cloud-узлах. Для SDN-контроллера именно \(AoS_{CTRL}\) является основной freshness-метрикой.

\section{Интерфейс отчетов и PHY-компоненты штрафа}

Минимальный пакет KPI публикуется как сочетание классов и опциональных оцененных численных значений. Для model checking существенны классы; численные значения могут использоваться внешним SDN-оптимизатором:
\begin{verbatim}
{
  SINRClass, BLERClass, CQIClass, IClass, DopplerClass,
  DelaySpreadClass, PowerClass, DRTClass, PilotDensityClass,
  PayloadSenseClass, PRSClass, BeamErrorClass, BlockageClass,
  PdClass, RfaClass, AccClass, CRBClass,
  BeamClass, MisClass, BMOverheadClass,
  AoSClass, CoverageClass, CapClass, ResourceShareClass,
  ChannelClass, SignalClass, SensingState, PHYState
}
\end{verbatim}

Если SDN-контроллеру требуется численный score, PHY может передавать нормированные компоненты нарушения контракта:
\begin{align}
\Pi_{SINR}&=\max\left(0,\frac{SINR_{min}-SINR_c}{SINR_{min}}\right),&
\Pi_{BLER}&=\max\left(0,\frac{BLER-BLER_{max}}{BLER_{max}}\right),\\
\Pi_{Pd}&=\max\left(0,\frac{Pd_{min}-Pd}{Pd_{min}}\right),&
\Pi_{Rfa}&=\max\left(0,\frac{Rfa-Rfa_{max}}{Rfa_{max}}\right),\\
\Pi_{Acc}&=\max\left(0,\frac{Acc_r-l_a}{l_a}\right)+\max\left(0,\frac{Acc_v-v_a}{v_a}\right),&
\Pi_{AoS}&=\max\left(0,\frac{AoS_{CTRL}-AoS_{max}}{AoS_{max}}\right),\\
\Pi_{Cap}&=\max\left(0,\frac{C_s^{min}-C_s}{C_s^{min}}\right).&&
\end{align}
Для beam management:
\begin{equation}
\Pi_{BM}=\max\left(0,\frac{\epsilon_b-\epsilon_{max}}{\epsilon_{max}}\right)+\max\left(0,\frac{p_{mis}-p_{mis}^{max}}{p_{mis}^{max}}\right)+\max\left(0,\frac{\Omega_{BM}-\Omega_{BM}^{max}}{\Omega_{BM}^{max}}\right).
\end{equation}
Агрегированный score
\begin{equation}
\Pi_{PHY}=\sum_i w_i\Pi_i,
\qquad w_i\ge 0,
\qquad \sum_iw_i=1
\end{equation}
не является частью обычного timed automaton, если используется как накопленная стоимость. Его следует трактовать либо как внешний SDN-score, вычисляемый после публикации отчета, либо как расширение модели до priced/weighted timed automata. В базовой TA-модели проверяются классы нарушений и bounded реакции, а не оптимальность \(\Pi_{PHY}\).

\section{Верификационные свойства и observer-автоматы}

В UPPAAL-подобной нотации можно задавать reachability и safety-свойства напрямую, но bounded response следует выражать observer-автоматами. Unbounded leads-to \(p\;--\!>\;q\) говорит только, что из \(p\) когда-нибудь следует \(q\); он не задает срок \(D\). В этой спецификации deadline трактуется как включительный: событие на границе \(c=D\) еще допустимо, а нарушение возникает только при \(c>D\). Поэтому observer-локации ожидания не имеют invariant \(c\le D\); иначе состояние \(c>D\) было бы недостижимо. Невмешательство observers обеспечивается тем, что наблюдаемые report/outcome-события объявлены как \code{broadcast chan}.

Базовые свойства:
\begin{verbatim}
A[] not deadlock
E<> A_PH.PHYNormal
E<> A_PH.PHYSensingDegraded
E<> A_PH.PHYCommunicationDegraded
E<> A_PH.PHYJointDegraded
A[] (A_PH.PHYNormal imply (comm_ok && sensing_qos_ok))
\end{verbatim}

Observer для bounded report после sensing degradation:
\begin{verbatim}
locations: Idle, WaitReport, Violation
clock: c_obs

Idle -- sensing_degraded? / c_obs := 0 --> WaitReport
WaitReport -- phy_kpi_report? guard c_obs <= D_report --> Idle
WaitReport -- c_obs > D_report --> Violation

query: A[] not ObsSenseReport.Violation
\end{verbatim}

Observer для freshness:
\begin{verbatim}
Idle -- aos_ctrl_expired? / c_obs := 0 --> WaitFreshnessLimited
WaitFreshnessLimited -- sensing_report? --> Idle
  guard: c_obs <= D_sense && SensingState == FreshnessLimited
WaitFreshnessLimited -- c_obs > D_sense --> Violation

query: A[] not ObsFreshness.Violation
\end{verbatim}

Observer для beam recovery:
\begin{verbatim}
Idle -- recovery_start? / c_obs := 0 --> WaitBeamOutcome
WaitBeamOutcome -- beam_restored? guard c_obs <= D_BM --> Idle
WaitBeamOutcome -- handover_hint? guard c_obs <= D_BM --> Idle
WaitBeamOutcome -- beam_failure? guard c_obs <= D_BM --> Idle
WaitBeamOutcome -- c_obs > D_BM --> Violation

query: A[] not ObsBeamRecovery.Violation
\end{verbatim}
С учетом исправленного перехода \(c_{rec}=D_{BM}\) в \code{BeamFailed} также допустимо проверять локальный invariant:
\begin{verbatim}
A[] (A_BM.BeamRecover imply c_rec <= D_BM)
\end{verbatim}
но это не заменяет проверку того, что recovery действительно завершается одним из разрешенных исходов.

Детерминированность классификаций проверяется через счетчики enabled-guards:
\begin{verbatim}
A[] (ch_enabled_count <= 1)
A[] (sq_enabled_count <= 1)
\end{verbatim}
или доказывается синтаксически через приоритетные функции \(\prio_{CH}\) и \(\prio_{SQ}\).

Согласованность assume--guarantee контрактов проверяется как отсутствие violation при выполнении конкретных boolean-предикатов. В UPPAAL-модели эти предикаты задаются как функции над состояниями и классами, например \code{ass\_ch()}, \code{gar\_ch()}, \code{ass\_sq()} и \code{ass\_ph()}:
\begin{verbatim}
bool ass_ch()  = ENV_CH.TimingOk && classes_in_range_ch;
bool gar_ch()  = A_CH.Reported && ChannelClass in CHANNEL_ENUM;
bool ass_sig() = ENV_MAC.ConfigAdmissible && classes_in_range_sig;
bool gar_sig() = A_SIG.Reported && SignalClass in SIGNAL_ENUM;
bool ass_bm()  = ENV_MAC.BeamCmdAdmissible && classes_in_range_bm;
bool gar_bm()  = A_BM.Reported && BeamClass in BEAM_ENUM;
bool ass_sq()  = fresh_child_reports && gar_ch() && gar_sig() && gar_bm();
bool gar_sq()  = A_SQ.Reported && SensingState in SENSING_ENUM;
bool ass_ph()  = gar_ch() && gar_sig() && gar_bm() && gar_sq();

A[] (ass_ch()  imply not A_CH.ContractViolation_CH)
A[] (ass_sig() imply not A_SIG.ContractViolation_SIG)
A[] (ass_bm()  imply not A_BM.ContractViolation_BM)
A[] (ass_sq()  imply not A_SQ.ContractViolation_SQ)
A[] (ass_ph()  imply not A_PH.ContractViolation_PH)
\end{verbatim}

\section{Ограничения применимости}

Предложенная формализация намеренно отделяет physical estimation от automata logic. Estimator или радиосимулятор оценивает \(SINR\), \(BLER\), \(Pd\), \(Rfa\), \(CRB\), \(p_{mis}\) и \(\Omega_{BM}\); timed automata получают конечные классы и проверяют своевременную реакцию. Если требуется доказывать распределения задержек, частоты отказов, среднюю стоимость или вероятность достижения состояния, модель должна быть расширена до probabilistic, stochastic или priced timed automata.

Качество результата зависит от soundness функции \(\alpha_{PHY}\). Пороговые значения должны выбираться из link budget, measurement confidence intervals, требований SLA или калибровочного симулятора. Грубая дискретизация может давать ложные counterexamples; слишком оптимистичная дискретизация может нарушить сохранение safety-свойств. Поэтому все границы классов в safety-критичных условиях относятся к худшему классу, а найденные counterexamples проверяются на физическую реализуемость вне model checker.

\section{Заключение}

PHY-уровень в SDN-управляемой ISAC-сети целесообразно задавать как композицию аппроксимированных контрактных timed automata. В строгой редакции модель содержит явную функцию \(\alpha_{PHY}\), конечные enum-классы для continuous и probabilistic KPI, assume--guarantee семантику, синхронизационные каналы, clocks, guards, resets, invariants, приоритеты классификации и observer-автоматы для deadlines.

Состояния автоматов отражают не детальную радиофизику, а проверяемые классы поведения: состояние канала, сигнальную конфигурацию, состояние луча, качество зондирования и агрегированное состояние PHY. Такая постановка честно ограничивает область применимости timed automata и одновременно делает модель полезной для формальной проверки SDN-управляемого поведения 6G ISAC-сети: MAC/SDN получает не произвольную таксономию PHY-состояний, а синхронизированную контрактную модель с проверяемыми сроками реакции.

\begin{thebibliography}{9}

\bibitem{li2026rethinking}
Y. Li, Y. Zhang, C. Masouros, S. Pollin, and F. Liu, ``Rethinking Signaling Design for ISAC: From Pilot-Based to Payload-Based Sensing,'' \emph{IEEE Communications Standards Magazine}, in press, 2026, doi: 10.1109/mcomstd.2025.3645941.

\bibitem{liu2024senscap}
G. Liu et al., ``SensCAP: A Systematic Sensing Capability Performance Metric for 6G ISAC,'' \emph{IEEE Internet of Things Journal}, vol. 11, no. 18, pp. 29438--29454, 2024, doi: 10.1109/JIOT.2024.3430502.

\bibitem{liu2026cooperative}
G. Liu et al., ``Cooperative Sensing for 6G ISAC: Concept, Key Technologies, Performance Evaluation, and Field Trial,'' \emph{Engineering}, vol. 56, pp. 130--148, 2026.

\bibitem{huang2026beam}
Y. Huang, J. Xu, M. A. Badiu, G. Chen, J. Coon, and M.-S. Alouini, ``ISAC-Enabled Low-Overhead Beam Management: Performance Analysis and Pilot Optimization,'' \emph{IEEE Transactions on Wireless Communications}, vol. 25, pp. 10702--10715, 2026, doi: 10.1109/TWC.2026.3653956.

\bibitem{alur1994timed}
R. Alur and D. L. Dill, ``A Theory of Timed Automata,'' \emph{Theoretical Computer Science}, vol. 126, no. 2, pp. 183--235, 1994.

\bibitem{uppaal2026queries}
UPPAAL Documentation, ``Query Semantics: Symbolic Queries,'' accessed 2026. Available: \url{https://docs.uppaal.org/language-reference/query-semantics/symb_queries/}.

\bibitem{anand2007afdx}
M. Anand and S. Vestal, ``Formal Modeling and Analysis of the AFDX Frame Management Design,'' in \emph{Proceedings of the IEEE International Conference on Engineering of Complex Computer Systems}, 2007.

\end{thebibliography}

\end{document}
